\documentclass[12 pt, letterpaper]{hmcpset}
\usepackage{hw-preamble}

\name{Jayden Thadani}
\class{MATH 142 --- Differential Geometry}
\assignment{Homework assignment \#8}
\duedate{29th March 2024}

\begin{document}

\begin{problem}[Problem 1.]
	\begin{enumerate}
		 \item Say $F, G : \mathbb{R}^{3} \to \mathbb{R}^{3}$ are isometries. Show that $\operatorname{sgn}(F \circ G) = \operatorname{sgn}(F) \operatorname{sgn}(G)$.
		\item Show that $F^{-1} : \mathbb{R}^{3} \to \mathbb{R}^{3}$ and $F$ have the same sign. That is, $\operatorname{sgn}(F) = \operatorname{sgn}(F^{-1})$.
	\end{enumerate}
\end{problem}

\begin{solution}
	\begin{enumerate}
		\item Suppose $F = T \circ C$ and $G = S \circ D$ where $T, S$ are translations $T : \mathbf{p} \mapsto \mathbf{p} + \mathbf{a}$ and $S : \mathbf{p} \mapsto \mathbf{p}$ and $C, D$ are orthogonal transformations. \\\\
			Then since $C$ is a linear map (all orthogonal transformations are linear), we have
			\[
				\begin{split}
					(F \circ G) (\mathbf{p}) & = F((S \circ D)(\mathbf{p})) \\
								 & = F(D(\mathbf{p}) + \mathbf{b}) \\
								 & = (T \circ C)(D(\mathbf{p}) + \mathbf{b}) \\
								 & = C(D(\mathbf{p}) + \mathbf{b}) + \mathbf{a} \\
								 & = (C \circ D)(\mathbf{p}) + (C(\mathbf{b}) + \mathbf{a}).
				\end{split}
			\]
			Thus, we can write $F \circ G$ as the composition $R \circ E$ where $E = C \circ D$ is an orthogonal transformation, and $R$ is translation by $C(\mathbf{b}) + \mathbf{a}$. Thus,
			\[
				\begin{split}
					\operatorname{sgn}(F \circ G) & = \det E \\
								      & = \det (C \circ D) \\
								      & = (\det C) (\det D) \\
								      & = \operatorname{sgn} (F) \operatorname{sgn} (G)
				\end{split}
			\]
			as desired, where the third equality is due to the multiplicity of the determinant for linear operators.
		\item We know that $F \circ F^{-1} = I$ where $I$ is the identity map on $\mathbb{R}^{3}$. Since $I$ is an orthogonal map itself, we have
			\[
				\begin{split}
					\operatorname{sgn} I & = \det I \\
							     & = 1.
				\end{split}
			\]
			Thus,
			\[
				\operatorname{sgn}(F) \operatorname{sgn}(F^{-1}) = 1.
			\] \\
			Since their product is positive, they must have the sign. Since we already know that the sign of any isometry has magnitude $1$, they must be equal.
	\end{enumerate}
\end{solution}

\pagebreak

\begin{problem}[Problem 2.]
	Recall that any orthogonal transformation $C : \mathbb{R}^{3} \to \mathbb{R}^{3}$ can be uniquely expressed in the form $C (\mathbf{x}) = A \mathbf{x}$ for some $3 \times 3$ orthogonal matrix $A$; we define $\operatorname{sgn}(C) = \det A$. A \textit{rotation} is an is an orthogonal transformation $C$  such that $\operatorname{sgn}(C) = 1$ .
	\begin{enumerate}
		\item Show that $A$ has eigenvalue $1$ as well as a unit eigenvector $\mathbf{e}_{3}$ satisfying $A \mathbf{e}_{3} = \mathbf{e}_{3}$.
		\item Prove that $C$ does, in fact, rotate $\mathbb{R}^{3}$ about an axis $L$. Explicitly, given a rotation $C$, show that there exists a number $\vartheta$ and points $\mathbf{e}_{1}, \mathbf{e}_{2}, \mathbf{e}_{3} \in \mathbb{R}^{3}$ with
			\[
				\mathbf{e}_{i} \cdot \mathbf{e}_{j} = \begin{cases}
					1 & \text{if } i = j \text{, and} \\
					0 & \text{otherwise}
				\end{cases}
			\]
			such that
			\[
				\begin{split}
					C(\mathbf{e}_{1}) & = \cos{\vartheta} \, \mathbf{e}_{1} + \sin{\vartheta} \, \mathbf{e}_{2} \\
					C(\mathbf{e}_{2}) & = - \sin{\vartheta} \, \mathbf{e}_{1} + \cos{\vartheta} \, \mathbf{e}_{2} \\
					C(\mathbf{e}_{3}) & = \mathbf{e}_{3}.
				\end{split}
			\]
	\end{enumerate}
\end{problem}

\pagebreak

\begin{problem}[Problem 3.]
	Let $\mathbf{a}$  be a point of $\mathbb{R}^{3}$  such that $\lVert \mathbf{a} \rVert = 1$ .
	\begin{enumerate}
		\item Prove that the formula
			\[
				C(\mathbf{p}) = \mathbf{a} \times \mathbf{p} + (\mathbf{p} \cdot \mathbf{a}) \mathbf{a}
			\]
			defines an orthogonal transformation $C : \mathbb{R}^{3} \to \mathbb{R}^{3}$.
		\item Describe its general effect on $\mathbb{R}^{3}$.
	\end{enumerate}
\end{problem}

\begin{solution}
	\begin{enumerate}
		\item First we will show that $C$ is linear.
			\[
				\begin{split}
					C(\alpha \mathbf{p} + \beta \mathbf{q}) & = \mathbf{a} \times (\alpha \mathbf{p} + \beta \mathbf{q}) + ((\alpha \mathbf{p} + \beta \mathbf{q}) \cdot \mathbf{a}) \mathbf{a} \\
								   & = \mathbf{a} \times (\alpha \mathbf{p}) + \mathbf{a} \times (\beta \mathbf{q}) + ((\alpha \mathbf{p}) \cdot \mathbf{a} +  (\beta \mathbf{q}) \cdot \mathbf{a}) \mathbf{a} \\
								   & = \alpha(\mathbf{a} \times \mathbf{p}) + \beta (\mathbf{a} \times \mathbf{q}) + \alpha(\mathbf{p} \cdot \mathbf{a}) \mathbf{a} + \beta (\mathbf{p} \cdot \mathbf{a}) \mathbf{a} \\
								   & = \alpha(\mathbf{a} \times \mathbf{p} + (\mathbf{p} \cdot \mathbf{a}) \mathbf{a}) + \beta (\mathbf{a} \times \mathbf{q} + (\mathbf{q} \cdot \mathbf{a}) \mathbf{a}) \\
								   & = \alpha C(\mathbf{p}) + \beta C(\mathbf{q}).
				\end{split}
			\] \\
			We will now show that $C$ preserves norms.
			\[
				\begin{split}
					\lVert C(\mathbf{p}) \rVert^{2} & = \lVert \mathbf{a} \times \mathbf{p} + (\mathbf{p} \cdot \mathbf{a}) \mathbf{a} \rVert^{2} \\
									& = \lVert \mathbf{a} \times \mathbf{p} \rVert^{2} + \lVert (\mathbf{p} \cdot \mathbf{a}) \mathbf{a} \rVert^{2} \\
									& = \lVert \mathbf{a} \rVert^{2} \lVert \mathbf{p} \rVert^{2} - (\mathbf{p} \cdot \mathbf{a})^{2} + (\mathbf{p} \cdot \mathbf{a})^{2} \lVert \mathbf{a} \rVert^{2} \\
									& = \lVert \mathbf{p} \rVert^{2}.
				\end{split}
			\]
			Here, the third step follows by Lagrange's identity, and the last step follows by the given fact that $\lVert \mathbf{a} \rVert = 1$. Since both are positive, we can take the square root on both sides to get $\lVert C(\mathbf{p}) \rVert = \lVert \mathbf{p} \rVert$. \\\\
			Using both of these, we can see that $C$ preserves distances. Specifically, 
			\[
				\begin{split}
					d(C(\mathbf{p}), C(\mathbf{q})) & = \lVert C(\mathbf{p}) - C(\mathbf{q}) \rVert \\
									& = \lVert C(\mathbf{p} - \mathbf{q}) \rVert \\
									& = \lVert \mathbf{p} - \mathbf{q} \rVert \\
									& = d(\mathbf{p}, \mathbf{q}).
				\end{split}
			\] \\
			Since $C$ preserves norms and distances, it is an orthogonal transformation.
		\item $C$ rotates any point by an angle of $\pi/2$ about the axis defined by $\operatorname{span}(\mathbf{a})$.
	\end{enumerate}
\end{solution}

\pagebreak

\begin{problem}[Problem 4.]
	Let $F : \mathbb{R}^{3} \to \mathbb{R}^{3}$ be an isometry of $\mathbb{R}^{3}$, and write $F = T \circ C$ for some translation $T$ and some orthogonal transformation $C$. Let $\beta : I \to \mathbb{R}^{3}$ be a unit speed curve in $\mathbb{R}^{3}$. Recall that the Frenet-Serret apparatus for $\overline{\beta} = F \circ \beta$ is
	\begin{align*}
		\overline{\mathbf{T}} & = F_{*} \circ \mathbf{T} & \overline{\mathbf{N}} & = F_{*} \circ \mathbf{N} & \overline{\mathbf{B}} & = \operatorname{sgn} F (F_{*} \circ \mathbf{B}) \\
		\overline{v} & = 1 & \overline{\kappa} & = \kappa & \overline{\tau} & = \operatorname{sgn} F \, \tau.
	\end{align*}
	\begin{enumerate}
		\item Show that if $\beta$ is a cylindrical helix, then $\overline{\beta} = F \circ \beta$ is also a cylindrical helix.
		\item The \textit{spherical image} $\sigma : I \to \mathbb{R}^{3}$ of $\beta$ is that function with the same Euclidean coordinates as that of $\beta'(t)$. That is, $\mathbf{T}(t) = [\sigma(t), \beta(t)]$. Show that, if $\beta$ has spherical image $\sigma$, then $\overline{\beta}$ has spherical image $C \circ \sigma$.
	\end{enumerate}
\end{problem}

\pagebreak

\begin{problem}[Problem 5.]
	Let $F$ be an isometry of $\mathbb{R}^{3}$. For each vector field $\mathbf{V}$ which sends a point $\mathbf{p}$ to the tangent vector $\mathbf{V}(\mathbf{p}) = [\mathbf{v}, \mathbf{p}]$, let $\overline{\mathbf{V}}$ be that vector field which sends a point $\mathbf{q} = F(\mathbf{p})$ to the tangent vector $F_{*}(\mathbf{V}(\mathbf{p}))$. Prove that \textit{isometries preserve covariant derivatives}. That is, show $\overline{\nabla_{\mathbf{V}} \mathbf{W}} = \nabla_{\overline{\mathbf{V}}} \overline{\mathbf{W}}$.
\end{problem}

\begin{solution}
	Let $F = T \circ C$ where $C$ is an orthogonal transformation and $T$ is a translation. Let $\mathbf{W}(\mathbf{p}) = [\mathbf{w}, \mathbf{p}]$ where $\mathbf{w}$ is a function $\mathbf{w} : \mathbb{R}^{3} \to \mathbb{R}^{3}$.
	We know, by definition, that $\nabla_{\mathbf{V}} \mathbf{W}(\mathbf{p}) = [D\mathbf{w} \Big |_{\mathbf{p}}(\mathbf{v}), \mathbf{p}]$. Also, notice that for $\mathbf{q} = F(\mathbf{p})$, we have $\overline{\mathbf{V}}(\mathbf{q}) = [(C \circ \mathbf{v} \circ F^{-1})(\mathbf{q}), \mathbf{q}]$, and similarly for $\overline{\mathbf{W}}$. Thus, for some $\mathbf{q} = F(\mathbf{p})$,
	\[
		\begin{split}
			\nabla_{\overline{\mathbf{V}}} \overline{\mathbf{W}}(\mathbf{q}) & = [D(C \circ \mathbf{w} \circ F^{-1}) \Big |_{\mathbf{p}} (C \circ \mathbf{v})(\mathbf{p}), \mathbf{q}] \\
									      & = \left [ DC \Big |_{\mathbf{w}(\mathbf{p})} \left ( D \mathbf{w} \Big |_{\mathbf{p}} \left ( DF^{-1} \Big |_{\mathbf{q}} (C(\mathbf{v}(\mathbf{p}))) \right ) \right ), \mathbf{q} \right ] \\
									      & = \left [ C \left ( D\mathbf{w} \Big |_{\mathbf{p}} (C^{-1}(C(\mathbf{v}(\mathbf{p})))) \right ), \mathbf{q} \right ] \\
									      & = \left [ C \left ( D \mathbf{w} \Big |_{\mathbf{p}} (\mathbf{v}(\mathbf{p})) \right ), \mathbf{q} \right ] \\
									      & = F_{*} \left ( \left [ D\mathbf{w} \Big |_{\mathbf{p}} (\mathbf{v}(\mathbf{p})), \mathbf{p} \right ] \right ) \\
									      & = F_{*}(\nabla_{\mathbf{V}} \mathbf{W}(\mathbf{p})) \\
									      & = \overline{\nabla_{\mathbf{V}} \mathbf{W}}{(\mathbf{q})}.
		\end{split}
	\]
	Here the second step follows by the multivariable chain rule, the third step follows by noticing that the orthogonal part of $F^{-1}$ is $C^{-1}$ and that the derivative of any linear map is itself, and the steps after those just follow by the definitions. \\\\
	Note that we used the notation $Df \Big |_{\mathbf{a}}$ to denote the linear map that is the total derivative of the function $f$, at the point $\mathbf{a}$, and the notation $Df \Big |_{\mathbf{a}}(\mathbf{v})$ to denote evaluating that linear map on the vector $\mathbf{v}$.
\end{solution}

\end{document}
